\documentclass[12pt, a4paper]{ctexart}
\usepackage[margin=2cm]{geometry}
\usepackage{libertine}

\usepackage{titlesec}
\usepackage{zhnumber}
\titleformat*{\section}{\Large\bfseries\raggedright}
\renewcommand{\thesection}{\zhnum{section}}
\renewcommand{\thesubsection}{\arabic{subsection}}

\setcounter{secnumdepth}{4}
\renewcommand{\theparagraph}{\thesubsubsection.\arabic{paragraph}}
\renewcommand{\paragraph}[1]{
    \refstepcounter{paragraph}
    {\bfseries\theparagraph\quad#1}\par
    \vspace{2pt}
    \noindent
}

\usepackage{enumitem}
\setlist[enumerate]{itemsep=2pt, parsep=0pt, partopsep=0pt, topsep=2pt}
\setlist[itemize]{itemsep=2pt, parsep=0pt, partopsep=0pt, topsep=2pt}
\linespread{1.2}

\usepackage{minted}
\setminted{
    breaklines=true, fontsize=\small, frame=lines,
    numbers=none, style=bw, tabsize=4
}
\usepackage{xparse}
\NewDocumentCommand{\mil}{ O{text} m }{ %
    \mintinline{#1}{#2} %
}

\usepackage{graphicx}
\usepackage{siunitx}

\begin{document}

\pagestyle{plain}
\thispagestyle{empty}

\noindent
\begin{tabular*}{\textwidth}{l @{\extracolsep{\fill}} r @{\extracolsep{6pt}} l}
    \LARGE{\textbf{实验报告}} & 计算机系统基础 & \textit{Introduction to Computer Systems} \\
\end{tabular*}\\\\
\begin{tabular*}{\textwidth}{l l}
    \textbf{报告标题: } & \textbf{实验环境安装与测试} \\
    \textbf{学号: } & 19240212 \\
    \textbf{姓名: } & 华博文 \\
    \textbf{日期: } & 2026 年 3 月 29 日 \\
\end{tabular*}\\
\rule[2ex]{\textwidth}{2pt}

\section{实验环境}

\subsection{操作系统}

macOS 26.4 arm64，Darwin 25.4.0，Apple M5。

\subsection{编程工具}

NeoVim v0.11.6，Apple clang 21.0.0。

\subsection{其他工具}

\LaTeX{}。

\section{实验目的}

\begin{enumerate}
    \item 掌握安装本实验的环境并以一简单程序 \mil{hello.c} 进行测试；
    \item 熟悉掌握 \mil{gcc} 编译命令和 \mil{objdump} 反汇编命令。
\end{enumerate}

\section{实验内容}

\subsection{环境安装}

因为已经安装了 Xcode Command Line Tools，所以直接使用 \mil{gcc} 和 \mil{objdump} 进行测试。

\subsection{使用 \mil{gcc} 完成简单 C 语言程序的编译和执行}

准备简单 C 语言程序 \mil{hello.c}，源代码如下：

\inputminted{c}{../src/hello.c}

使用下列 gcc 命令将 hello.c 程序依次转换为相应的可执行文件。

\begin{minted}{bash}
gcc -E hello.c -o hello.i # 预处理
gcc -S hello.i -o hello.s # 编译
gcc -c hello.s -o hello.o # 汇编
gcc hello.o -o hello # 链接
gcc -O0 -g -no-pie -fno-pic hello.c -o hello2 # 直接编译
\end{minted}

分别执行可执行文件 \mil{hello} 和 \mil{hello2}，均输出 \mil{hello world!}。

\subsection{使用 \mil{objdump} 命令实现反汇编}

用 \mil{objdump} 命令对 \mil{hello.o} 和 \mil{hello} 反汇编，并保存到文本文件：

\begin{minted}{bash}
objdump -S hello.o > hello_o.asm
objdump -S hello > hello.asm
\end{minted}

比较 \mil{hello.o} 和 \mil{hello} 的反汇编内容差异：

\inputminted{text}{../res/hello_o.asm}

\inputminted{text}{../res/hello.asm}

用 \mil{objdump} 命令对 \mil{hello2} 反汇编：

\begin{minted}{bash}
objdump -S hello2 > hello2.asm
\end{minted}

比较 \mil{hello} 和 \mil{hello2} 的反汇编内容差异：

\inputminted{text}{../res/hello2.asm}

\section{结论}

本实验完成了简单 C 语言程序的编译与执行流程，熟悉了从源代码到可执行文件的各个阶段，包括预处理、编译、汇编和链接。同时，利用 \mil{objdump} 工具对生成的目标文件和可执行文件进行了反汇编，深入理解了不同编译选项对生成代码的影响。实验中，对比了 \mil{hello} 和 \mil{hello2} 的反汇编结果，加深了对编译器工作原理和程序底层实现的理解。

\end{document}
